\section{Evaluation plan}
\label{sec:eval-plan}

The main thing I need to show is that the system can take a real refactoring trace and produce explanations that a developer would actually find useful. Concretely, given (i) a repository, (ii) a chosen start and end revision, and (iii) an observed trace of intermediate snapshots, the system should (a) score the observed trajectory, (b) construct a plausible higher-scoring reference trajectory between the same endpoints, and (c) highlight a small set of divergence points with short ``what changed / why it matters / what could have been done instead'' explanations.

My primary plan is to collect refactoring episodes from humans using my custom IDE plugin, so I have traces in exactly the format the system is designed for. These episodes should include both IDE-supported refactorings and manual edits, and capture whatever signals I can log reliably (e.g., checkpoints/snapshots, refactoring commands, and build/test outcomes where available). This gives me realistic traces with the right granularity, without relying on commit history as the only source of ``steps''.

On top of these real episodes, I will create controlled variants by deliberately injecting known bad-process patterns (e.g., unnecessary detours, rework, churn spikes, temporary quality regressions, or broken intermediate builds/tests). The point is to create a rough ground truth for where the divergence \emph{should} be. I can then check whether the system finds these injected divergence points and whether the constructed reference trajectory avoids the injected detours while still reaching the same end state. I will measure this mainly by whether the reference reduces the ``bad'' signals I injected (e.g., effort/churn, temporary smell/metric spikes, intermediate build/test failures), rather than only by endpoint quality.

I am also considering mining additional episodes from open-source Java repositories, especially small and relatively self-contained refactoring PRs or multi-commit refactoring sequences, but I am not committing to that yet. If I do it, the goal would be to stress-test robustness to real-world messiness (mixed edits, inconsistent commit granularity, incomplete separation between refactoring and feature work) and to increase the diversity of codebases beyond what I can collect through the plugin alone.

I will also run a small human evaluation to checkt that the insights produced by the tool on how their refactoring process could have been improved are valuable.